% 例 3：非对称区间概率图
% 用法：% 例 3：非对称区间概率图
% 用法：% 例 3：非对称区间概率图
% 用法：% 例 3：非对称区间概率图
% 用法：\input{assets/tikz/R032_example_asymmetric_interval.tex}

\begin{center}
	\begin{tikzpicture}[
		x=1.52cm,
		y=4.6cm,
		>=Stealth,
		line cap=round,
		line join=round,
		every node/.style={font=\small},
		axis/.style={-Stealth, black!60, line width=0.7pt}
	]
		\path[use as bounding box] (-3.85,-0.24) rectangle (3.85,0.55);
		\fill[blue!24]
			(-1,0) -- plot[domain=-1:2, samples=90] (\x,{0.4*exp(-\x*\x/2)}) -- (2,0) -- cycle;
		\draw[axis] (-3.65,0) -- (3.66,0) node[right] {$z$};
		\draw[blue!65!black, line width=1.05pt]
			plot[domain=-3.55:3.55, samples=150] (\x,{0.4*exp(-\x*\x/2)});
		\foreach \x/\lab in {-1/{-1},0/{0},2/{2}}{
			\draw[black!55] (\x,-0.012)--(\x,0.012);
			\node[below=1pt] at (\x,-0.004) {$\lab$};
		}
		\draw[dashed, blue!70!black, line width=0.8pt] (-1,0)--(-1,0.26);
		\draw[dashed, blue!70!black, line width=0.8pt] (2,0)--(2,0.09);
		\node[font=\bfseries\large, text=blue!70!black] at (0,0.49)
			{非对称区间：两个端点分别标准化};
		\node[
			draw=blue!65!black,
			fill=white,
			rounded corners=5pt,
			inner xsep=8pt,
			inner ysep=3pt,
			text=blue!70!black
		] at (0,0.23) {$\Phi(2)-\Phi(-1)\approx0.8186$};
		\node[font=\scriptsize, text=black!65] at (0,-0.18)
			{$(-1,8)\xrightarrow{\ z=(x-2)/3\ }(-1,2)$};
	\end{tikzpicture}
\end{center}


\begin{center}
	\begin{tikzpicture}[
		x=1.52cm,
		y=4.6cm,
		>=Stealth,
		line cap=round,
		line join=round,
		every node/.style={font=\small},
		axis/.style={-Stealth, black!60, line width=0.7pt}
	]
		\path[use as bounding box] (-3.85,-0.24) rectangle (3.85,0.55);
		\fill[blue!24]
			(-1,0) -- plot[domain=-1:2, samples=90] (\x,{0.4*exp(-\x*\x/2)}) -- (2,0) -- cycle;
		\draw[axis] (-3.65,0) -- (3.66,0) node[right] {$z$};
		\draw[blue!65!black, line width=1.05pt]
			plot[domain=-3.55:3.55, samples=150] (\x,{0.4*exp(-\x*\x/2)});
		\foreach \x/\lab in {-1/{-1},0/{0},2/{2}}{
			\draw[black!55] (\x,-0.012)--(\x,0.012);
			\node[below=1pt] at (\x,-0.004) {$\lab$};
		}
		\draw[dashed, blue!70!black, line width=0.8pt] (-1,0)--(-1,0.26);
		\draw[dashed, blue!70!black, line width=0.8pt] (2,0)--(2,0.09);
		\node[font=\bfseries\large, text=blue!70!black] at (0,0.49)
			{非对称区间：两个端点分别标准化};
		\node[
			draw=blue!65!black,
			fill=white,
			rounded corners=5pt,
			inner xsep=8pt,
			inner ysep=3pt,
			text=blue!70!black
		] at (0,0.23) {$\Phi(2)-\Phi(-1)\approx0.8186$};
		\node[font=\scriptsize, text=black!65] at (0,-0.18)
			{$(-1,8)\xrightarrow{\ z=(x-2)/3\ }(-1,2)$};
	\end{tikzpicture}
\end{center}


\begin{center}
	\begin{tikzpicture}[
		x=1.52cm,
		y=4.6cm,
		>=Stealth,
		line cap=round,
		line join=round,
		every node/.style={font=\small},
		axis/.style={-Stealth, black!60, line width=0.7pt}
	]
		\path[use as bounding box] (-3.85,-0.24) rectangle (3.85,0.55);
		\fill[blue!24]
			(-1,0) -- plot[domain=-1:2, samples=90] (\x,{0.4*exp(-\x*\x/2)}) -- (2,0) -- cycle;
		\draw[axis] (-3.65,0) -- (3.66,0) node[right] {$z$};
		\draw[blue!65!black, line width=1.05pt]
			plot[domain=-3.55:3.55, samples=150] (\x,{0.4*exp(-\x*\x/2)});
		\foreach \x/\lab in {-1/{-1},0/{0},2/{2}}{
			\draw[black!55] (\x,-0.012)--(\x,0.012);
			\node[below=1pt] at (\x,-0.004) {$\lab$};
		}
		\draw[dashed, blue!70!black, line width=0.8pt] (-1,0)--(-1,0.26);
		\draw[dashed, blue!70!black, line width=0.8pt] (2,0)--(2,0.09);
		\node[font=\bfseries\large, text=blue!70!black] at (0,0.49)
			{非对称区间：两个端点分别标准化};
		\node[
			draw=blue!65!black,
			fill=white,
			rounded corners=5pt,
			inner xsep=8pt,
			inner ysep=3pt,
			text=blue!70!black
		] at (0,0.23) {$\Phi(2)-\Phi(-1)\approx0.8186$};
		\node[font=\scriptsize, text=black!65] at (0,-0.18)
			{$(-1,8)\xrightarrow{\ z=(x-2)/3\ }(-1,2)$};
	\end{tikzpicture}
\end{center}


\begin{center}
	\begin{tikzpicture}[
		x=1.52cm,
		y=4.6cm,
		>=Stealth,
		line cap=round,
		line join=round,
		every node/.style={font=\small},
		axis/.style={-Stealth, black!60, line width=0.7pt}
	]
		\path[use as bounding box] (-3.85,-0.24) rectangle (3.85,0.55);
		\fill[blue!24]
			(-1,0) -- plot[domain=-1:2, samples=90] (\x,{0.4*exp(-\x*\x/2)}) -- (2,0) -- cycle;
		\draw[axis] (-3.65,0) -- (3.66,0) node[right] {$z$};
		\draw[blue!65!black, line width=1.05pt]
			plot[domain=-3.55:3.55, samples=150] (\x,{0.4*exp(-\x*\x/2)});
		\foreach \x/\lab in {-1/{-1},0/{0},2/{2}}{
			\draw[black!55] (\x,-0.012)--(\x,0.012);
			\node[below=1pt] at (\x,-0.004) {$\lab$};
		}
		\draw[dashed, blue!70!black, line width=0.8pt] (-1,0)--(-1,0.26);
		\draw[dashed, blue!70!black, line width=0.8pt] (2,0)--(2,0.09);
		\node[font=\bfseries\large, text=blue!70!black] at (0,0.49)
			{非对称区间：两个端点分别标准化};
		\node[
			draw=blue!65!black,
			fill=white,
			rounded corners=5pt,
			inner xsep=8pt,
			inner ysep=3pt,
			text=blue!70!black
		] at (0,0.23) {$\Phi(2)-\Phi(-1)\approx0.8186$};
		\node[font=\scriptsize, text=black!65] at (0,-0.18)
			{$(-1,8)\xrightarrow{\ z=(x-2)/3\ }(-1,2)$};
	\end{tikzpicture}
\end{center}
